
\documentclass{report}

\input{../../preamble}
\input{../../macros}
\input{../../letterfonts}

\title{Assignment 06}
\author{Ryan Lin}
\date{}

\begin{document}
\maketitle
\section{Question 1:}

Let $ x = A_{1}u + b_1 $, $ y = A_{2}v +b_{2} $. 

Original Optimization: 
\begin{align*}
\min_{u,v}: \quad & \lVert (A_1u + b_1) - (A_2v + b_2)\rVert_{2} \\
\text{s.t}: \quad & \lVert u \rVert_{2}  \le 1 \\
		  & \lVert v \rVert_{2} \le 1
\end{align*}


But this is not conic, which has the form $ \lVert Az + b\rVert_{2} \le c^{T}z +d $. 


Use a variable $ t $ 
\begin{align*}
\min_{u,v,t}: \quad & t \\
\text{s.t}: \quad & \lVert (A_1u + b_1) - (A_2v +b_2) \rVert_{2} \le t \\
		  & \lVert u \rVert_{2} \le 1 \\
		  & \lVert v \rVert_{2} \le 1
\end{align*}


We can show that this is a SOCP by defining $ z = \begin{bmatrix}
	u \\
	v\\
	t
\end{bmatrix} $ 
\[
\lVert (A_1u + b_1) - (A_2v +b_2) \rVert_{2} = \lVert \begin{bmatrix}
	A_1 &-A_2 & 0
\end{bmatrix} \cdot  z + \underbrace{b_1 - b_2}_{d}\rVert_{2} \le \begin{bmatrix}
	0 &0 &1
\end{bmatrix} \cdot z
\]

Constraint 2: 
\[
	\lVert \begin{bmatrix}
		I_{2} &0 &0
	\end{bmatrix} \cdot z \le 1
\]
Constraint 3: 
\[
	\lVert \begin{bmatrix}
		0 &I_2 &0
	\end{bmatrix} \cdot z \le 1
\]


\newpage

\subsection{a}
Let $ 1 = w^{T}w $, $ x=t, y= a^{T}x + b$

Thus, 
\begin{align*}
	t(a^{T}x + b) \ge 1 &= xy \ge w^{T}w \\
\end{align*}

\begin{align*}
	xy \ge w^{T}w \\
	\left\lVert \begin{bmatrix}
		2w\\
		x-y
	\end{bmatrix} \right \rVert_{2} \le x + y\\
	\left\lVert \begin{bmatrix}
		2\\ 
		t - a^{T}x - b
	\end{bmatrix} \right\rVert \ge t + (a^{T}x + b)
\end{align*}

From this, we can see that 

\begin{align*}
		x &=  \begin{bmatrix}
		x\\
		t
	\end{bmatrix}\\
A &= \begin{bmatrix}
		0 & 0\\
		-a^{T} & 1
	\end{bmatrix}\\
b&= \begin{bmatrix}
	2\\
	-b
\end{bmatrix} \\
c &= \begin{bmatrix}
	a^{T} \\1
\end{bmatrix} \\
	d &= b
\end{align*}

\subsection{b}

We can use the auxiliary variable $ t $ to help put this in standard form. 

We add the constraints: 
\begin{align*}
	t_{i} &\ge \frac{1}{a_{i}^{T}x + b_{i}}\\
	t_{i}(a_{i}^{T}x + b_{i}) \ge 1	
\end{align*}

This is the same form as part a. 

\begin{align*}
\min_{x, t}: \quad & \sum_{i=1}^{p} t_{i} \\
\text{s.t}: \quad & t_{i}(a_{i}^{T}x + b_{i}) \ge 1, i=1,\dots,p\\
		  &a_{i}^{T} + b_{i} > 0, i=1, \dots, p\\
		  &c_{j}^{T}x + d_{j} \ge 0, j=1,\dots, q\\
		  &t_{i} \ge 0, i=1, \dots,p
\end{align*}

Which can be rewritten as: 
\begin{align*}
\min_{x, t}: \quad & \sum_{i=1}^{p} t_{i} \\
\text{s.t}: &\lVert Ax + b \rVert_{2} \le c^{T}x + b \\
		  &a_{i}^{T} + b_{i} > 0, i=1, \dots, p\\
		  &c_{j}^{T}x + d_{j} \ge 0, j=1,\dots, q\\
		  &t_{i} \ge 0, i=1, \dots,p
\end{align*}




\end{document}
